\documentclass[12pt]{article}
\usepackage[utf8]{inputenc}
\usepackage[spanish]{babel}
\usepackage{amsmath}
\usepackage{amsfonts}
\usepackage{amssymb}
\usepackage{geometry}
\geometry{letterpaper, margin=2.5cm}

\begin{document}

\begin{center}
    \Large \textbf{Evaluación de Examen 1: Unidad I} \\
    \large Física para Ciencias de la Salud (005-1713) \\
    \vspace{0.5cm}
    \textbf{Estudiante:} Michelle López \quad \textbf{C.I.:} 33.483.425
    \vspace{0.3cm} \\
    \large \textbf{Calificación Final: 10/10 puntos}
\end{center}

\vspace{0.5cm}

\section*{Análisis de Resultados}

\subsection*{1. Ángulo plano (Conversión a radianes)}
\textbf{Procedimiento y resultado:} Plantea correctamente la conversión multiplicando directamente por el factor $\frac{\pi}{180^\circ}$. Evalúa la fracción $\frac{75\pi}{180}$ obteniendo como resultado un valor decimal válido y redondeado de $0,42\pi$ Rad. \\
\textbf{Veredicto:} \textbf{Correcto} (2/2 pts).

\subsection*{2. Sistemas de coordenadas (Longitud del hueso)}
\textbf{Procedimiento y resultado:} Excelente resolución analítica. Emplea la fórmula de la distancia entre dos puntos y sustituye las coordenadas de manera precisa. Desarrolla aritméticamente los binomios al cuadrado y la suma ($\sqrt{36+64} = \sqrt{100}$) para concluir con el resultado exacto de $10$ cm. \\
\textbf{Veredicto:} \textbf{Correcto} (2/2 pts).

\subsection*{3. Vectores (Fuerza resultante)}
\textbf{Procedimiento y resultado:} La estudiante evidencia un gran orden al separar el cálculo por componentes ($\vec{F}_{Rx}$ y $\vec{F}_{Ry}$). Efectúa las sumas algebraicas respetando la ley de los signos para obtener 30 y 60 respectivamente. Finalmente, ensambla el vector resultante de forma impecable: $(30\hat{i} + 60\hat{j})$ N. \\
\textbf{Veredicto:} \textbf{Correcto} (2/2 pts).

\subsection*{4. Potencias de diez (Notación científica y prefijo)}
\textbf{Procedimiento y resultado:} Transforma la cifra decimal correctamente a notación de potencias de base diez ($12,5 \cdot 10^{-6}$ m). Adicionalmente, aunque con un trazo muy tenue de lápiz, identifica el exponente $-6$ con su correspondiente símbolo de prefijo del S.I. ($\mu$), escribiendo $\mu$m a un lado de su respuesta. \\
\textbf{Veredicto:} \textbf{Correcto} (2/2 pts).

\subsection*{5. Sistemas de unidades (Conversión de densidad)}
\textbf{Procedimiento y resultado:} Magistral aplicación de los factores de conversión en cadena. Expresa las relaciones fraccionarias correctas para masa y volumen, simplificando de forma clara la ecuación ($1,03 \cdot 1000$). Obtiene un resultado numérico y de unidades totalmente exacto: $1030$ Kg/m$^3$. \\
\textbf{Veredicto:} \textbf{Correcto} (2/2 pts).

\vspace{0.5cm}
\section*{Comentarios Generales y Sugerencias}
¡Felicidades por obtener la calificación máxima! El examen destaca por su excelente presentación, orden procedimental en los desarrollos algebraicos y un entendimiento profundo de la notación científica y las conversiones de unidades.
\begin{itemize}
    \item Como única y pequeña sugerencia, se recomienda remarcar un poco más los resultados secundarios (como el símbolo $\mu$m en la Pregunta 4) para asegurar su fácil visibilidad durante la lectura.
\end{itemize}

\end{document}